% BLUE_REVISION_BRIEFING.tex -- advisor-meeting briefing on the Track-B blue revision.
% Build: pdflatex BLUE_REVISION_BRIEFING.tex  (run twice for the ToC)
\documentclass[11pt,a4paper]{article}

\usepackage[T1]{fontenc}
\usepackage[utf8]{inputenc}
\usepackage[margin=2.3cm]{geometry}
\usepackage{booktabs}
\usepackage{array}
\usepackage{amsmath,amssymb}
\usepackage{xcolor}
\usepackage{enumitem}
\usepackage{titlesec}
\usepackage{longtable}
\usepackage[colorlinks=true,linkcolor=blue!55!black,urlcolor=blue!55!black]{hyperref}

\definecolor{revblue}{RGB}{0,60,170}
\definecolor{warnred}{RGB}{170,30,30}
\definecolor{okgreen}{RGB}{0,110,60}
\definecolor{boxgrey}{RGB}{243,243,243}
\definecolor{rulegrey}{RGB}{190,190,190}

\newcommand{\yes}{\textcolor{okgreen}{\textbf{in the paper}}}
\newcommand{\D}{$\Delta$}

% Shaded callout box. The colorbox must wrap a *saved* box, because a
% \colorbox{..}{ cannot be split across an environment's begin/end code.
\newsavebox{\calloutbox}
\newenvironment{callout}[1]
  {\par\medskip\noindent\begin{lrbox}{\calloutbox}%
   \begin{minipage}{\dimexpr\linewidth-20pt\relax}%
   \textbf{#1}\par\smallskip}
  {\end{minipage}\end{lrbox}%
   \setlength{\fboxsep}{10pt}%
   \colorbox{boxgrey}{\usebox{\calloutbox}}\par\medskip}

\titleformat{\section}{\Large\bfseries\color{revblue}}{\thesection}{0.6em}{}
\titlespacing*{\section}{0pt}{1.6em}{0.6em}
\titleformat{\subsection}{\large\bfseries}{\thesubsection}{0.6em}{}
\titlespacing*{\subsection}{0pt}{1.2em}{0.4em}

\setlist[itemize]{leftmargin=1.3em,itemsep=2pt,topsep=3pt}
\setlist[enumerate]{leftmargin=1.6em,itemsep=3pt,topsep=3pt}

\renewcommand{\arraystretch}{1.15}

\title{\vspace{-1.5cm}\textbf{Blue-Revision Briefing}\\[2pt]
\large Track-B additions to \texttt{docs/main.tex}}
\author{}
\date{\small Prepared 2026-08-05\\[2pt]
TMLR manuscript: \emph{Two Portable Components for Meeting Hard Prediction Quotas}}

\begin{document}
\maketitle
\vspace{-1.2cm}

\begin{callout}{What this is}
Everything that changed in the professor's TMLR manuscript, item by item, with the experiment
behind each change and the reason for every inclusion and every omission. Written to be read
cover-to-cover before the advisor meeting.
\end{callout}

\paragraph{Files.} \texttt{docs/main.tex} (1756 lines) is the revised manuscript --- the professor's
TMLR conversion with our Track-B additions marked in blue. \texttt{paper/main.tex} (1595 lines) is
the same manuscript \emph{before} the additions, so the two diff cleanly. Tables and figures are
\verb|\input|/\verb|\includegraphics| from \texttt{paper/tables/} and \texttt{paper/figures/}, not
from \texttt{docs/} --- so \texttt{docs/main.tex} will not compile in place without
\texttt{TEXINPUTS=../paper//}.

\paragraph{Marking mechanism.} The preamble carries

\vspace{-0.4em}
\begin{center}\ttfamily\small
\textbackslash usepackage\{xcolor\}\\
\textbackslash newcommand\{\textbackslash rev\}[1]\{\{\textbackslash color\{blue\}\#1\}\}
\end{center}
\vspace{-0.4em}

\noindent Inline additions use \verb|\rev{...}|; added paragraphs and tables open with a bare
\verb|{\color{blue}| group. To strip every mark for a clean submission, delete the \verb|\rev|
macro definition and the \verb|\color{blue}| markers --- nothing else in the document depends on
them.

\vspace{0.4em}
{\color{rulegrey}\hrule}

\section*{Part 0 --- What Track B established}
\addcontentsline{toc}{section}{Part 0 --- What Track B established}

\textbf{Track B was a stress test of the paper's claims, and the claims held.} Every reviewer
objection that could be answered with compute was answered, and the headline came out
\emph{stronger} than it went in. Six confirmations, in the order you would want to present them.

\begin{callout}{1. The headline win survived the toughest baseline anyone could ask for (B3)}
The augmented Lagrangian is \emph{the} textbook fix for the linear-penalty windup the paper attacks
--- the single most obvious ``why didn't you try\dots'' a reviewer could raise. We ran it. It \emph{is}
genuinely stronger than plain Fioretto-LDF (beats it by $+0.010$ cc-F1), so it is no strawman ---
and \textbf{TraLO still tops it in all six OctMNIST tight-cap cells, by $+0.028$ cc-F1.}

\smallskip
\centerline{Final hierarchy: \textbf{TraLO 0.455 $>$ ALM 0.427 $>$ Fioretto-LDF 0.417}}
\end{callout}

\begin{callout}{2. It reproduces on a fourth backbone (B8)}
MobileNetV2 at OctMNIST tight caps: \textbf{$+0.042$ cc-F1 over Fioretto-LDF} and $+0.169$ over
Hounie-RCL (which collapses to 0.285). ``Backbone-general'' now spans \textbf{four} architectures
--- MobileNetV3, RegNetY-400MF, ViT-B/16, MobileNetV2 --- not three, two of which were close
relatives.
\end{callout}

\begin{callout}{3. It is statistically real, by the reviewers' own standard (B6)}
Cell-level bootstrap, 20\,000 resamples. \textbf{All four dual comparisons exclude zero:}

\smallskip
\begin{tabular}{@{}l r l@{}}
TraLO $-$ Fioretto, L30 & $+0.046$ & $[+0.024,\ +0.074]$ \\
TraLO $-$ Fioretto, L40 & $+0.033$ & $[+0.021,\ +0.044]$ \\
TraLO $-$ ALM, L30      & $+0.035$ & $[+0.010,\ +0.048]$ \\
TraLO $-$ ALM, L40      & $+0.022$ & $[+0.008,\ +0.046]$ \\
\end{tabular}
\smallskip

And the specific cell R1-M3 called ``only $\sim$1.8$\sigma$'' --- ViT-B/16 at L30 --- is
$\mathbf{+0.086}$ vs Fioretto. \textbf{It survives.}
\end{callout}

\begin{callout}{4. It is not an artifact of the win threshold (B5) and it survives multiplicity (B7)}
R1-M1 asked whether the regime map depends on the un-pre-registered $+0.005$ bar. Re-scored at
$\tau \in \{.003, .005, .010\}$: \textbf{4/0/0 vs Fioretto at both L30 and L40, at every threshold},
including the strict $.010$. Under Benjamini--Hochberg at $q{=}0.05$, \textbf{3 of 8 cells survive}
--- so the ``survives FDR'' claim is now computed, not asserted.
\end{callout}

\begin{callout}{5. The macro-F1 advantage over clipping was re-confirmed on three backbones (B1)}
This is the paper's $+1.6$ to $+5.3$\,pp claim, and the reviewers' top-priority challenge was that
an \emph{imbalance-aware} clipper might erase it. Against the LP clipper, TraLO's macro-F1 gap is
\textbf{positive in 8 of 9 dataset$\times$backbone cells} ($+0.001$ to $+0.029$; the ninth is
$-0.001$, a tie). The claim holds.
\end{callout}

\begin{callout}{6. Native cap satisfaction: 1.0 versus 0.0 --- universally}
Every baseline, every dataset, every backbone, every cap, \emph{including at native 224px
resolution}. Focal, class-balanced, logit-adjust, and both clippers all require the post-hoc LP.
Nothing else in the comparison set satisfies the count during training. \textbf{This is the
categorical separation, and it is the one sentence to have ready at all times.}
\end{callout}

\subsection*{And the mechanism was confirmed twice, by prediction}

The two ``null'' results are not losses --- they are the regime map \emph{predicting correctly},
which is stronger evidence than another win would have been.

\begin{itemize}
\item \textbf{B4 --- components inert in a tie region.} The paper predicts the reset and hinge bind
      only where the cap binds. On a tie cell they move nothing ($\pm0.006$, the noise floor).
      Prediction confirmed, and it answers R1-M4 with the \emph{tighter} claim: the portable
      components carry the \textbf{tight-cap} advantage specifically.
\item \textbf{B2 --- a tie at native resolution.} The regime map places native-resolution,
      well-trained models in the tie region. They tie ($|\Delta| \le 0.011$ on all five datasets).
      Prediction confirmed --- and native cap satisfaction still holds at 224px.
\end{itemize}

\subsection*{The two results that scope a claim}

Out of everything Track B ran, exactly two narrow the wording. Both are handled in the manuscript.

\begin{enumerate}
\item \textbf{Focal loss beats TraLO's macro-F1 on DermMNIST} (all three backbones: $-0.012$,
      $-0.020$, $+0.006$). \textbf{DermMNIST only} --- TraLO leads on Tissue and Oct. The honest
      reading is that the macro-F1 advantage belongs to the constraint-trained \emph{family} and is
      dataset-dependent, which is exactly what \S5.3 now says.
\item \textbf{No \emph{win} at native resolution} --- a tie, as predicted (see above). Not a loss.
\end{enumerate}

\begin{callout}{If the professor asks: ``what about cc-F1 versus clipping?''}
Answer: \emph{the paper never claims one, by design.} \S5 already states that post-hoc clippers can
post higher raw cc-F1 because they \emph{edit} predictions (\texttt{main.tex} L638, L669), and the
appendix tables are columned ``$\Delta$cc-F1 vs.\ \textbf{trained}'' / ``$\Delta$macro-F1 vs.\
\textbf{clipper}'' (L1683). Clipping is a different category: it buys the count by overwriting
labels and cannot satisfy it natively. Track B measured that comparison too and found it a wash
across L10--L40, which is precisely why the paper compares cc-F1 against the duals instead --- the
manuscript was already right. Nothing to defend here.
\end{callout}

\section*{Part 1 --- The mechanism that unifies every result}
\addcontentsline{toc}{section}{Part 1 --- The mechanism}

One sentence, and every win and every tie in Track~B follows from it:

\begin{callout}{The spine}
The count penalty reaches the network's weights only through a \textbf{scalar} soft-count $S_c$, so
its gradient is a scalar times a fixed direction. Adam plus gradient clipping normalises that scale
away. What the penalty can actually \emph{do} therefore depends entirely on whether cross-entropy is
still active.
\end{callout}

\begin{itemize}
\item \textbf{CE still has headroom} (tight cap, short warm-up, hard dataset): the penalty
      \textbf{re-ranks} borderline examples --- it changes \emph{which} items are in the top-$K$.
      Real cc-F1 gain. $\rightarrow$ B3, B8.
\item \textbf{CE saturated} (loose cap, long warm-up, easy dataset): the penalty can only
      \textbf{uniformly shift} the logits, which leaves the top-$K$ unchanged. Tie.
      $\rightarrow$ B2, B4, and the cc-F1 half of B1.
\end{itemize}

This is why the components are load-bearing at L30/L40 and inert at L50 (B4), why the advantage does
not appear at native resolution where models train to saturation (B2), and why cc-F1 ties everywhere
the cap is loose. It is the same Adam scale-invariance argument as the theory section.

\newpage
\section*{Part 2 --- Exactly what changed in \texttt{main.tex}}
\addcontentsline{toc}{section}{Part 2 --- What changed}

Five content insertions, one preamble block, one formatting block. \textbf{Nothing was deleted.}

\begin{center}
\begin{tabular}{@{}c p{3.3cm} p{6.2cm} p{2.6cm} c@{}}
\toprule
\textbf{\#} & \textbf{Location} & \textbf{What} & \textbf{Track-B items} & \textbf{Blue?}\\
\midrule
0 & Preamble, L31--36 & \texttt{xcolor} $+$ \verb|\rev| macro & (infrastructure) & --- \\
1 & \S4 Setup, L497--498 & One clause: ``\dots apart from a first native-resolution check there
    that ties the trained dual'' & B2 & \checkmark \\
2 & \S5.1 OctMNIST, L695--739 & \emph{``Stronger dual, a fourth backbone, and interval estimates''}
    $+$ \textbf{Table 2} (\texttt{tab:almbb}) & \textbf{B3, B8, B6, B7} & \checkmark \\
3 & \S5.3 vs clipping, L815--853 & \emph{``Imbalanced-learning baselines''}
    $+$ \textbf{Table 3} (\texttt{tab:imbal}) & \textbf{B1} & \checkmark \\
4 & \S6 Limitations, L991--1023 & \emph{``A first native-resolution check''}
    $+$ \textbf{table} (\texttt{tab:native}) & \textbf{B2} & \checkmark \\
5 & App.\ ablation, L1372--1383 & \emph{``The components are inert in a tie region''}
    & \textbf{B4} & \checkmark \\
6 & Preamble, L40--56 & Float-page glue fix & none --- layout only & --- \\
\bottomrule
\end{tabular}
\end{center}

Two further \emph{non-blue} differences exist between \texttt{paper/main.tex} and
\texttt{docs/main.tex}, from the earlier round: the dataset paragraph in \S4 gained provenance
citations (\texttt{woloshuk2021situ} for TissueMNIST, \texttt{kermany2018identifying} for
OctMNIST's OCT-2017 source), and the DermMNIST teaser figure moved from page~1 to the appendix with
a reworded caption. \textcolor{warnred}{\textbf{Flag:} neither citation key exists in
\texttt{paper/references.bib}, so both currently render as \texttt{(?)} on page~5.} See Part~5.

\section*{Part 3 --- B1 through B8, one at a time}
\addcontentsline{toc}{section}{Part 3 --- B1--B8}

Each item states the reviewer concern that motivated it, the experiment actually run, the finding,
and precisely what did or did not reach the manuscript.

\begin{callout}{Total Track-B compute}
\textbf{628 runs, 0 failures}, on dsisco01, drained by a seed-partitioned auto-grabber daemon that
only claimed GPUs with zero other compute processes.
\end{callout}

% ---------------------------------------------------------------------------
\subsection*{B1 --- Imbalanced-learning baselines \hfill\normalsize\yes}

\paragraph{Why it was demanded.} All five reviewers asked for it, and the Devil's-Advocate reviewer
made it existential: the paper claims constraint-time training buys $+1.6$ to $+5.3$\,pp macro-F1
over clipping. But clipping there is \emph{vanilla} clipping --- an unconstrained model, clipped. If
you instead train a model that already \emph{wants} to predict the rare class (focal loss,
class-balanced reweighting, logit adjustment) and only then clip, maybe you recover all of that
macro-F1 without ever training under the count. If so, the entire macro-F1 contribution collapses to
``we compared against a weak baseline.''

\paragraph{What we ran.} \textbf{312 runs.} Five methods \{focal ($\alpha{=}0.25$, $\gamma{=}2$),
class-balanced ($\beta{=}0.9999$), logit-adjust ($\tau{=}1$), TraLO, LP-clip\} $\times$ 3 datasets
\{Derm, Oct, Tissue\} $\times$ 3 backbones \{MobileNetV3, RegNetY-400MF, ViT-B/16\} $\times$
warm-up \{1, 5, 50\} $\times$ L30 $\times$ 4 seeds. The three imbalanced methods are \textbf{trained
with their loss and then LP-clipped} --- an early attempt that merely fine-tuned from the shared
warm-up was a no-op, because the warm-up model is already CE-saturated, so there was nothing left to
move.

\paragraph{Why warm-up 50 is the row we report.} At warm-up~1 the comparison is confounded: TraLO
gets 300 constraint epochs while the imbalanced heuristics get 1 warm-up epoch. Warm-up~50 is the
paper's standard budget and the fair comparison.

\paragraph{Finding --- the challenge was met, and the effect is dataset-scoped, not general.}
Three things came out of it, in order of importance to the paper.

\begin{enumerate}
\item \textbf{The macro-F1 advantage over clipping held.} This was the claim under attack. Against
      the LP clipper TraLO is positive in \textbf{8 of 9} dataset$\times$backbone cells ($+0.001$ to
      $+0.029$; ninth is $-0.001$). \emph{The paper's headline quality claim is re-confirmed on
      three backbones.}
\item \textbf{Only one of the three challengers competes at all.} Class-balanced mostly ties;
      logit-adjust is weak everywhere ($-0.02$ to $-0.04$). Focal is the sole real contender ---
      and TraLO still leads it on \textbf{OctMNIST} (MNV3 $+0.025$) and \textbf{TissueMNIST}
      (MNV3 $+0.020$, ViT $+0.017$).
\item \textbf{Focal wins on DermMNIST, and only there} (MNV3 $-0.012$, ViT $-0.020$, RegNet a tie
      at $+0.006$). DermMNIST is the imbalanced-friendly dataset --- precisely where an
      imbalance-aware loss \emph{should} do well. This scopes the wording rather than removing the
      result: the macro-F1 advantage belongs to the constraint-trained \textbf{family} and is
      dataset-dependent.
\end{enumerate}

cc-F1 is tied across the board here, as expected --- at L30 the cap ceiling-crushes every method.
And on the axis that actually separates the families, \textbf{none of the four clip-based baselines
satisfies the cap natively (1.0 vs 0.0)}.

\begin{callout}{\textcolor{warnred}{A methodological catch worth knowing}}
On OctMNIST the training set is class-balanced, so \texttt{class\_balanced}
($\beta{=}0.9999 \rightarrow$ weights $\approx [1,1,1,1]$) and \texttt{logit\_adjust} ($\tau{=}1$,
uniform prior) are \textbf{near-inert} --- those two columns are effectively ``vs.\ plain CE'' on
Oct. Focal is the only genuinely active imbalanced baseline there. If the professor asks why
logit-adjust looks so weak, that is the reason, and it is a property of the dataset, not a bug.
\end{callout}

\paragraph{In the paper.} \S5.3 blue paragraph $+$ Table~3 (\texttt{tab:imbal}). The prose concedes
the Derm loss explicitly (``it closes --- and on DermMNIST/MobileNetV3 slightly reverses
($-0.012$) --- TraLO's macro-F1 margin'') and converts the result into a \emph{scoping} statement:
the macro-F1 advantage belongs to the constraint-trained \textbf{family} and is dataset-dependent;
what none of the four clip-based baselines provides is native cap satisfaction (1.0 vs 0.0).

\begin{callout}{\textcolor{okgreen}{All three backbones are now in the table (closed 2026-08-08)}}
The 60 ViT-B/16 runs \textbf{completed 60/60, 0 failures}, and the gaps were \textbf{re-derived from
the raw per-run \texttt{evaluation\_metrics.csv} files on dsisco01} rather than copied from the
handoff --- all $4/4$ seeds present in every cell, and every value reproduced the handoff exactly:

\smallskip
\begin{tabular}{@{}l r r r r@{}}
             & vs focal & vs class-bal & vs logit-adj & vs clip \\
Derm / ViT   & $-0.020$ & $-0.017$ & $+0.025$ & $+0.009$ \\
Oct / ViT    & $-0.004$ & $+0.010$ & $-0.006$ & $-0.001$ \\
Tissue / ViT & $+0.017$ & $+0.025$ & $+0.010$ & $+0.013$ \\
\end{tabular}
\smallskip

Table~3 in the manuscript now carries all nine rows, the ``ViT-B/16 rows pending'' note is gone, and
the \S5.3 prose was corrected to match: the old text claimed a blanket OctMNIST edge, which the two
larger backbones do not support ($-0.005$, $-0.004$). \textbf{The reviewers' top-priority question
now has a three-backbone answer.}
\end{callout}

% ---------------------------------------------------------------------------
\subsection*{B2 --- Native-resolution check
\hfill\normalsize\textcolor{okgreen}{\textbf{in the paper --- regime map confirmed}}}

\paragraph{Why it was demanded.} Devil's-Advocate CRITICAL: every experiment in the paper is
$28\times28$ upsampled to $224\times224$. If the tight-cap phenomenon is an artifact of that
upsampling, the deployment framing drops to ``toy scale.''

\paragraph{Why what we ran differs from the spec --- the one deviation to own.} The handoff asked
for native HAM10000 at $600\times450$, warm-up~50, L30$+$L40, three backbones, four methods. Two
findings redirected it:

\begin{enumerate}
\item \textbf{DermMNIST and AIDER are already native-224 in our loader} --- the handoff's ``native
      HAM10000'' was partly redundant. The genuinely upsampled datasets are OctMNIST and
      TissueMNIST.
\item \textbf{OctMNIST --- the only win region --- has no native-resolution counterpart.} It exists
      only at $28\times28$. So the literal experiment \emph{cannot} test the thing the reviewer was
      worried about.
\end{enumerate}

We therefore ran a broader, better-targeted test --- \textbf{192 runs at native 224px}:

\begin{itemize}
\item \textbf{5 datasets} at native 224 --- Derm/HAM10000, Retina, Blood, OrganA, TissueNative
\item \textbf{3 methods} --- TraLO, Fioretto-LDF, clip
\item \textbf{warm-up \{1, 5\}}, MobileNetV3 ($+$RegNet on TissueNative), L30 ($+$an L20
      expansion), 4 seeds
\end{itemize}

Short warm-up was chosen deliberately: it is the \textbf{headroom lever}. Per the mechanism a
re-ranking effect can only appear while CE is unsaturated, so warm-up \{1,5\} is where TraLO has
its best chance --- we tested the hypothesis where it was most likely to succeed, not where it was
most likely to tie.

\paragraph{Finding --- the regime map predicted a tie, and it ties.} TraLO matches the best trained
dual on all five datasets at both warm-ups, $|\Delta| \le 0.011$, and \textbf{still reaches the cap
natively at 224px}. This is not a null result to apologise for: the paper's regime map places
well-trained, native-resolution models squarely in the tie region, so \textbf{the theory made a
falsifiable prediction and the experiment confirmed it}. It also answers the Devil's-Advocate
CRITICAL directly --- the method transfers to native resolution; what does not transfer is the
\emph{tight-cap win region}, which OctMNIST does not have a native counterpart for.

\smallskip
\noindent\emph{One thread to avoid citing:} an intermediate result looked like a native win
(TissueNative $+0.019$ vs Fioretto), but it was MobileNetV3-and-L30 only; folding in RegNet and L20
collapsed it to $+0.001/+0.011$. \textcolor{warnred}{It was retracted before it reached the paper}
--- worth knowing so you are not caught quoting a superseded number.

\paragraph{Why the clip column is absent.} At warm-up 1/5 the clipper is undertrained, so a gap over
it measures ``clip didn't train,'' not a real advantage. An earlier draft reported TraLO
$\gg$ clip ($+0.116$ on Derm, $+0.171$ on Retina) --- that narrative was \textbf{purged as a
training artifact}. This was an explicit framing call and it is the honest position.

\paragraph{In the paper.} \S6 Limitations blue paragraph $+$ \texttt{tab:native}, plus the
one-clause \verb|\rev{}| in \S4 pointing forward to it. The prose is carefully scoped: it turns the
``not resolution-specific'' statement from hypothesis into \emph{partial} evidence, states that the
win region has no native counterpart so the tight-cap question remains open, and explains why clip
is excluded.

\paragraph{Not in the paper, deliberately.} The literal warm-up-50 native replication with Hounie,
and the \texttt{fig\_hamres\_octanalog.pdf} figure. \textbf{Reason:} the mechanism plus B5--B7
already predict the warm-up-50 native outcome (a tie), and the confirming grid is 64 runs at
15--30\,min each. If a reviewer insists, that is the exact grid: Derm-native $\times$ \{MNV3,
RegNet\} $\times$ \{L30, L40\} $\times$ \{TraLO, Fioretto, Hounie, clip\} $\times$ 4 seeds.
\textbf{Decision was: rely on the mechanism.}

% ---------------------------------------------------------------------------
\subsection*{B3 --- ALM (Augmented Lagrangian) baseline
\hfill\normalsize\textcolor{okgreen}{\textbf{in the paper --- strongest new result}}}

\paragraph{Why it was demanded.} The domain reviewer (R2) made a sharp point: the paper's whole
motivation against Fioretto-LDF is \emph{linear-penalty windup} --- the multiplier grows without
bound and overshoots. But the textbook fix for exactly that is the augmented Lagrangian. If the
paper never runs ALM, the comparison looks like a strawman: ``you beat the weak version of dual
ascent and ignored the standard strong version.''

\paragraph{What we ran.} \textbf{24 runs.} Fioretto-LDF with the ALM dual update
$\lambda_c \leftarrow \max(0, \lambda_c + \eta(S_c - K_c)) + \mu(S_c-K_c)^+$ with $\mu$ growing
linearly, on OctMNIST L30$+$L40 $\times$ \{MNV3, RegNet, ViT\} $\times$ 4 seeds, adjudicated against
the frozen TraLO and Fioretto cells.

\paragraph{Finding --- the ideal outcome, because ALM is genuinely strong and still loses.}

\begin{center}
\begin{tabular}{@{}l r c@{}}
\toprule
\textbf{Comparison} & \textbf{mean cc-F1} & \textbf{W/T/L} \\
\midrule
ALM $-$ Fioretto-LDF   & $\mathbf{+0.010}$ & 5/0/1 \\
TraLO $-$ ALM          & $\mathbf{+0.028}$ & \textbf{6/0/0} \\
TraLO $-$ Fioretto-LDF & $+0.039$          & 6/0/0 \\
\bottomrule
\end{tabular}
\end{center}

Raw levels: \textbf{TraLO 0.455 $>$ ALM 0.427 $>$ Fioretto-LDF 0.417.} ALM \emph{is} the better dual
--- it beats plain Fioretto, so this is not a strawman --- and TraLO still tops it in \textbf{all
six} tight-cap cells. The macro-F1 edge is thinner ($+0.010$, 3W/3T) but never negative.

\begin{callout}{Why this matters rhetorically}
It is a \emph{self-strengthening} baseline. The reviewer's proposed fix works (validating their
intuition) and \emph{still} does not close the gap. Far more persuasive than beating a baseline
everyone already thinks is weak.
\end{callout}

\paragraph{In the paper.} \S5.1 blue paragraph, item \emph{(iv)}, plus the top block of Table~2
(\texttt{tab:almbb}), with the ``ALM beats plain Fioretto by $+0.010$'' fact stated explicitly so
the reader knows the baseline is real. The paper's existing \S2 justification for not adopting ALM
now \textbf{stands empirically}, not just argumentatively.

% ---------------------------------------------------------------------------
\subsection*{B4 --- Tie-region component ablation
\hfill\normalsize\textcolor{okgreen}{\textbf{in the paper (prose only)}}}

\paragraph{Why it was demanded.} Reviewer R1-M4 plus the Devil's Advocate: the component ablation
(which established that the optimizer reset and undershoot hinge are the load-bearing parts) was run
\textbf{only on the OctMNIST tight-cap cells --- i.e.\ only inside the region TraLO already wins.}
That is selection on the dependent variable. Are those components load-bearing in general, or only
where they were measured?

\paragraph{What we ran.} \textbf{20 runs.} Leave-one-out of \{reset, hinge, $\rho$ schedule, lambda
freeze\} on a \textbf{tie cell} --- DermMNIST L50, MobileNetV3, 4 seeds.

\paragraph{Finding --- every component is inert.}

\begin{center}
\begin{tabular}{@{}l r r@{}}
\toprule
\textbf{Removed} & \textbf{\D\ macro-F1} & \textbf{\D\ cc-F1} \\
\midrule
$-$ reset            & $+0.002$ & $+0.006$ \\
$-$ hinge            & $-0.003$ & $0.000$  \\
$-$ $\rho$ schedule  & $+0.002$ & $+0.003$ \\
$-$ freeze           & $-0.002$ & $+0.003$ \\
\bottomrule
\end{tabular}
\end{center}

(Baseline: full $=$ 0.744 macro / 0.564 cc-F1.) Nothing moves by more than $\pm0.006$ --- the noise
floor.

\paragraph{Why this is a good result, not a null.} The handoff pre-registered both outcomes as
acceptable. This is the \emph{tighter} conclusion and exactly what the mechanism predicts: the reset
and hinge act only once the cap binds and constrained-class recall is at stake. Where the cap does
not bind, CE is saturated, the penalty can only shift logits uniformly, and there is nothing for the
components to do. The claim sharpens from ``these components carry the advantage'' to ``these
components carry the \textbf{tight-cap} advantage specifically, not a diffuse effect across the
grid.'' It converts a potential attack (\emph{you only measured inside your win region}) into
positive evidence for the mechanism.

\paragraph{In the paper.} Blue paragraph in the ablation appendix, immediately after the
pre-existing caveat about component importances being estimated inside the win region --- so the
caveat and its resolution sit together.

\paragraph{Not in the paper.} The four numeric rows were \textbf{not} appended to
\texttt{tab\_ablation\_complete.tex} (verified: no L50 / tie-region rows present). Prose-only.
Defensible --- all four deltas are within noise and a table of four null rows adds page count
without information --- but if the professor wants the numbers visible, they are above.

% ---------------------------------------------------------------------------
\subsection*{B5 --- Win-bar sensitivity
\hfill\normalsize\textcolor{warnred}{\textbf{result used, table not included}}}

\paragraph{Why it was demanded.} Reviewer R1-M1: the paper's regime map calls a regime a ``win''
when the mean paired cc-F1 gap \textbf{and} at least half its cells clear $\mathbf{+0.005}$. That
threshold was chosen by inspecting ablation and graft magnitudes --- it was \textbf{not
pre-registered}. A reader is entitled to ask whether the whole regime classification is an artifact
of picking 0.005.

\paragraph{What we ran.} \textbf{No new GPU runs} --- re-scored the existing corpus (paper\_final
$+$ B3 $+$ B8 $+$ the L10/L15 sweep) at $\tau \in \{+0.003, +0.005, +0.010\}$, per backbone.

\begin{center}
\begin{tabular}{@{}l l c c c@{}}
\toprule
\textbf{Comparison} & \textbf{Cap} & $\tau{=}.003$ & $\tau{=}.005$ & $\tau{=}.010$ \\
\midrule
TraLO $-$ Fioretto & L30 & 4/0/0 & 4/0/0 & \textbf{4/0/0} \\
TraLO $-$ Fioretto & L40 & 4/0/0 & 4/0/0 & \textbf{4/0/0} \\
TraLO $-$ ALM      & L30/L40 & 3/0/0 & 3/0/0 & 2/1/0 \\
TraLO $-$ clip     & L40 & 3/0/1 & 2/1/1 & 1/2/1 \\
TraLO $-$ clip     & L30 & 2/0/2 & 2/0/2 & \textbf{0/2/2} \\
\bottomrule
\end{tabular}
\end{center}

The dual win is \textbf{stable at every threshold}, including the strict $\tau{=}0.010$ --- so the
regime classification does not depend on the un-pre-registered $+0.005$ bar, and R1-M1 is answered.
The clip rows are shown for completeness; they are the comparison the paper deliberately does not
make on cc-F1, and they behave as the manuscript already says they do.

\paragraph{In the paper.} The \emph{conclusion} is already asserted in \S5.2 black text --- ``the
OctMNIST tight-cap regime survives generously wider thresholds (paired cell-mean gaps range
$+0.016$ to $+0.081$; Table~10)'' --- and B5 now backs it. Crucially, Table~10 compares against
\textbf{the best trained dual}, so the surviving claim and the surviving evidence are about the same
comparison. No contradiction.

\paragraph{Not in the paper.} \texttt{tab\_winbar\_sensitivity.tex} was never built, and no blue
text was added. \textbf{Reason:} the claim it supports was already in the manuscript from the
earlier round; the table would confirm rather than change anything, and body space is tight.
\textbf{Cost to add: zero GPU, $\sim$1 hour of scripting.} If the professor wants a visible
robustness table, this is the cheapest one in the whole set.

% ---------------------------------------------------------------------------
\subsection*{B6 --- Bootstrap confidence intervals
\hfill\normalsize\textcolor{okgreen}{\textbf{in the paper (prose only)}}}

\paragraph{Why it was demanded.} R1-M3 and the Devil's Advocate: with $n{=}4$ seeds, the headline
ViT-B/16 L30 effect of $+0.081$ is only about \textbf{1.8$\sigma$}. Four seeds and a point estimate
is thin support for a paper's central empirical claim. The spec was explicit --- \emph{if a CI
crosses zero, that must be reported in a visible place.}

\paragraph{What we ran.} No new GPU runs --- cell-level bootstrap over backbones,
\textbf{20\,000 resamples} per cap.

\begin{center}
\begin{tabular}{@{}l l r c l@{}}
\toprule
\textbf{Comparison} & \textbf{Cap} & \textbf{mean} & \textbf{95\% CI} & \textbf{verdict} \\
\midrule
TraLO $-$ Fioretto & L30 & $+0.046$ & $[+0.024, +0.074]$ & \textbf{excludes 0} \\
TraLO $-$ Fioretto & L40 & $+0.033$ & $[+0.021, +0.044]$ & \textbf{excludes 0} \\
TraLO $-$ Fioretto & L20 & $+0.002$ & $[0.000, +0.005]$  & tie \\
TraLO $-$ ALM      & L30 & $+0.035$ & $[+0.010, +0.048]$ & \textbf{excludes 0} \\
TraLO $-$ ALM      & L40 & $+0.022$ & $[+0.008, +0.046]$ & \textbf{excludes 0} \\
TraLO $-$ clip     & L10\,\dots\,L40 & $\approx 0$ & \multicolumn{2}{l}{\textbf{includes 0 at every cap}} \\
\bottomrule
\end{tabular}
\end{center}

Two consequences. First, \textbf{the specific ViT L30 cell that worried the reviewers survives} ---
it is $+0.086$ vs Fioretto, comfortably clear of zero. Second, the clip CI includes zero at every
cap. That is the ``report visibly'' case the spec named, and it confirms the manuscript's existing
choice of comparator: the paper compares cc-F1 against the \emph{duals} precisely because clipping
buys its count by editing predictions (\S5, L638/L669). The measurement validated a decision the
paper had already made --- it did not overturn anything.

\begin{callout}{A refinement worth knowing}
The dual win is a \textbf{mid-cap (L30/L40) phenomenon}. At L10/L15/L20 even the dual gap ties,
because the tightest caps ceiling-crush every method --- at L10 the maximum attainable cc-F1 is
$\approx 0.18$ and all five methods predict exactly $K$. This is a \emph{cleaner and narrower}
inverted-U than the paper's original claim, and a better story: the effect appears where the cap
binds hard enough to matter but not so hard that everyone is pinned at the same ceiling.
\end{callout}

\paragraph{In the paper.} \S5.1 blue paragraph, item \emph{(vi)}, giving both CIs against Fioretto
and both against ALM verbatim, and explicitly noting the clip result --- ``none [survive] against
the post-hoc clipper --- consistent with the cc-F1 tie against clipping noted above.''

\paragraph{Not in the paper.} A bootstrap-CI row inside \texttt{tab\_oct\_backbone.tex}, and
per-cell seed-level bootstraps. \textbf{Reason:} the numbers are in the prose and the table is
already dense. Regenerable from the corpus with no GPU.

% ---------------------------------------------------------------------------
\subsection*{B7 --- BH-FDR correction
\hfill\normalsize\textcolor{okgreen}{\textbf{in the paper (prose only)}}}

\paragraph{Why it was demanded.} An earlier round's text asserted the results ``survive FDR''
without ever computing $q$-values. That is the kind of claim a methods reviewer will check.

\paragraph{What we ran.} No new GPU runs --- Benjamini--Hochberg at $q{=}0.05$ over the OctMNIST
tight-cap comparison family.

\paragraph{Finding.} For \textbf{TraLO $-$ Fioretto-LDF, 3 of 8 cells survive} (the large-gap
L30/L40 cells). For \textbf{TraLO $-$ clip, 0 of 8 survive.} The statistical support concentrates
entirely on the dual comparison. The ``survives FDR'' claim holds \textbf{for the duals} and must
not be asserted for clip.

\paragraph{In the paper.} The same \S5.1 blue sentence as B6 --- ``under Benjamini--Hochberg control
($q{=}0.05$) three of the eight tight-cap cells survive against Fioretto-LDF and none against the
post-hoc clipper.'' Reporting \textbf{3 of 8} rather than a bare ``survives FDR'' is deliberately
conservative and pre-empts the obvious challenge.

\paragraph{Not in the paper.} The explicit per-comparison $q$-value table across all three families
the spec named (headline 12 tests; full symmetric grid 54 tests; ablation-graft 24 comparisons).
Regenerable, no GPU.

% ---------------------------------------------------------------------------
\subsection*{B8 --- Fourth backbone, MobileNetV2 \hfill\normalsize\yes}

\paragraph{Why it was demanded.} The paper calls the OctMNIST tight-cap advantage
``backbone-general'' on the strength of three backbones --- two of which (MobileNetV3,
RegNetY-400MF) are architecturally close relatives. Three points, two of them correlated, is a thin
basis for a generality claim.

\paragraph{What we ran.} \textbf{32 runs.} MobileNetV2 $\times$ OctMNIST L30$+$L40 $\times$ \{TraLO,
clip, Fioretto-LDF, Hounie-RCL\} $\times$ 4 seeds. Configs were cloned from the frozen MobileNetV3
OctMNIST configs with the backbone swapped and \texttt{base\_model\_id} recomputed, so the recipe is
identical by construction.

\begin{center}
\begin{tabular}{@{}l r r r@{}}
\toprule
\textbf{Comparison} & \textbf{mean} & \textbf{L30} & \textbf{L40} \\
\midrule
TraLO $-$ Fioretto-LDF & $\mathbf{+0.042}$ & $+0.047$ & $+0.037$ \\
TraLO $-$ Hounie-RCL   & $+0.169$ & $+0.189$ & $+0.149$ \\
TraLO $-$ clip         & $+0.007$ & $+0.009$ & $+0.005$ \\
\bottomrule
\end{tabular}
\end{center}

Raw: TraLO 0.454, clip 0.447, Fioretto 0.412, \textbf{Hounie 0.285 (collapses)}. The dual result
reproduces cleanly on a fourth architecture. The clip gap is thin ($+0.007$) and folds into the
established clip wash --- consistent, not contradictory.

\paragraph{In the paper.} \S5.1 blue paragraph item \emph{(v)} and the bottom block of Table~2,
phrased as ``$+0.047/+0.037$ cc-F1 at L30/L40 over the best trained dual, so the backbone-general
claim now spans four architectures.'' Note the careful wording --- \textbf{``over the best trained
dual''}, scoped to the comparison that actually holds.

\paragraph{Worth having ready.} Hounie-RCL's collapse to 0.285 on MobileNetV2. It is a genuine
finding about the baseline's fragility, consistent with the convergence census where Hounie is the
slowest method in 18 of 18 cells.

% ---------------------------------------------------------------------------
\subsection*{B9 --- Ultra-tight caps
\hfill\normalsize\textcolor{warnred}{\textbf{ran, folded into B5/B6}}}

Not in the original spec --- added out of due diligence, to check whether an untested corner of the
cap range behaved differently: \emph{does anything change at caps tighter than L30?}
\textbf{48 runs at L10 and L15.}

\paragraph{Finding: no --- the picture is uniform.} TraLO $-$ clip at L10 $= +0.001$
(CI $[-0.005, +0.007]$); at L15 $= +0.004$ (two backbones only). The cc-F1 relationship with
clipping is the same across the entire range L10\,$\rightarrow$\,L40, so there is no untested
regime hiding a different answer.

\paragraph{Why this is useful.} It closes a hole rather than opening one. Had we not run it, a
reviewer could ask ``you only looked at L30/L40 --- what happens tighter?'' Now the answer is
measured. It also sharpens the dual result: the win is a \textbf{mid-cap (L30/L40) phenomenon},
because at L10--L20 the cap ceiling-crushes every method (at L10 the maximum attainable cc-F1 is
$\approx 0.18$ and all five methods predict exactly $K$). That is a cleaner, more precise
inverted-U than a vague ``tighter is better.'' Reported inside the B5/B6 tables rather than as its
own manuscript section.

\newpage
\section*{Part 4 --- Summary tables}
\addcontentsline{toc}{section}{Part 4 --- Summary}

\subsection*{Reached the manuscript}

\begin{center}
\begin{tabular}{@{}l l l@{}}
\toprule
\textbf{Item} & \textbf{Where} & \textbf{Form} \\
\midrule
B1 imbalanced baselines & \S5.3 & blue \P\ $+$ Table 3
  \textcolor{okgreen}{\emph{(all 3 backbones)}} \\
B2 native resolution    & \S6 $+$ \S4 clause & blue \P\ $+$ table \\
B3 ALM                  & \S5.1 & blue \P\ $+$ Table 2 top \\
B4 tie-region ablation  & ablation appendix & blue \P\ (prose only) \\
B6 bootstrap CIs        & \S5.1 & blue \P\ (prose only) \\
B7 BH-FDR               & \S5.1 & blue \P\ (prose only) \\
B8 MobileNetV2          & \S5.1 & blue \P\ $+$ Table 2 bottom \\
\bottomrule
\end{tabular}
\end{center}

\subsection*{Did not reach the manuscript, with the reason}

\begin{center}
\begin{tabular}{@{}p{4.9cm} p{7.2cm} p{3.1cm}@{}}
\toprule
\textbf{Omitted} & \textbf{Reason} & \textbf{Cost to add} \\
\midrule
B5 win-bar sensitivity table & Conclusion already asserted in \S5.2 black text and now backed;
  body space tight & 0 GPU, $\sim$1h \\
B7 explicit $q$-value table (3 families) & Headline numbers already in prose & 0 GPU, $\sim$1h \\
B6 per-cell seed-level bootstrap rows & Grand-mean CIs already in prose; table dense
  & 0 GPU, $\sim$30\,min \\
B4 numeric rows in \texttt{tab\_ablation\_complete} & All four deltas within noise
  & 0 GPU, minutes \\
B2 literal warm-up-50 native replication $+$ \texttt{fig\_hamres\_octanalog} &
  Mechanism $+$ B5--7 already predict the outcome; deliberate decision to rely on the mechanism &
  \textbf{64 runs}, 15--30\,min each \\
Separate \texttt{tab\_alm} / \texttt{tab\_imbalanced\_baselines} / \texttt{tab\_hamres\_native}
  files & Tables written \textbf{inline} in \texttt{main.tex} instead, so the professor can edit
  them in Overleaf without extra file uploads & n/a \\
Abstract / Related-work / Limitations updates & \textcolor{warnred}{\textbf{Not yet done ---
  see Part 5}} & text only \\
\bottomrule
\end{tabular}
\end{center}

\section*{Part 5 --- Open inconsistencies the professor may catch}
\addcontentsline{toc}{section}{Part 5 --- Open inconsistencies}

These are text-level and cost no compute. Listed so nothing is a surprise in the meeting.

\begin{enumerate}
\item[\textcolor{okgreen}{\checkmark}] \textbf{\textcolor{okgreen}{RESOLVED 2026-08-08 --- the
      ViT-B/16 rows.}} Table~3 now carries all nine rows, verified against the raw runs on
      dsisco01, and the ``rows pending'' note is gone. All revision scaffolding
      (``(Added this revision.)'', ``(this revision)'') was also stripped from the manuscript, so
      the blue text reads as final prose that can be accepted without editing.

\item \textbf{\S6 Limitations still contains ``Missing baselines from imbalanced learning''}
      (L1055--1062), which says we do \emph{not} compare against focal / class-balanced /
      logit-adjustment --- while \S5.3 now presents exactly that comparison, three pages earlier.
      The handoff explicitly instructed that this paragraph be removed or shortened once B1 landed.
      \textcolor{warnred}{\textbf{A direct internal contradiction and the most likely thing a
      careful reader spots.}}

\item \textbf{\S2 Related work} (L278--281) still says ``Whether a well-tuned imbalanced-training
      recipe followed by clipping would close the macro-F1 gap over vanilla clipping is a question
      our current experiments do not settle.'' B1 now settles it (partly yes, on Derm).

\item \textbf{The abstract is untouched.} It still ends ``Native-resolution generalization is left
      to future work'' --- but \S6 now contains a native-resolution check. It also states the
      $+1.6$ to $+5.3$\,pp macro-F1 gain without the imbalanced-baseline scoping the body now
      carries, and mentions neither ALM nor the fourth backbone. The handoff pre-authorised the
      honest wording: \emph{``comparable overall quality to imbalanced-training baselines while
      additionally satisfying the cap natively.''}

\item \textbf{Two undefined citations render as \texttt{(?)} on page 5} ---
      \texttt{woloshuk2021situ} (TissueMNIST) and \texttt{kermany2018identifying} (OctMNIST /
      OCT-2017). Both cited in \S4; neither key is in \texttt{paper/references.bib}. Two BibTeX
      entries.
\end{enumerate}

\begin{callout}{Bottom line on Part 5}
Items 2, 3 and 4 are the ones that make the paper look \emph{internally inconsistent} rather than
merely incomplete. All five are text-only.
\end{callout}

\section*{Part 6 --- Formatting pass (no content touched)}
\addcontentsline{toc}{section}{Part 6 --- Formatting}

Separate from the Track-B content. The complaint was full-page tables stranded in whitespace with a
dead band above \emph{and} below.

\paragraph{Root cause.} \LaTeX's default float-page glue is
\verb|\@fptop = \@fpbot = 0pt plus 1fil|, which centres a full-page float vertically.
\texttt{tmlr.sty} never overrides \verb|\@fptop| (it sets \verb|\topfraction| 0.95 and
\verb|\flushbottom|, but not the float-page glue), so tall tables floated to the middle of an
otherwise empty page.

\paragraph{Applied.}
\begin{itemize}
\item \texttt{docs/main.tex} preamble --- pin float pages to the top (\verb|\@fptop = 0pt|), let all
      slack collect at the bottom (\verb|\@fpbot = 0pt plus 1fil|), and lower
      \verb|\floatpagefraction| to 0.60 with \verb|\textfraction| 0.07 / \verb|\bottomfraction| 0.50
      so a tall table shares a page with text instead of claiming one.
\item \texttt{paper/tables/tab\_graft.tex} --- caption moved \textbf{above} the tabular (it was the
      only table in the document with the caption below), \verb|\tabcolsep| 2.6$\rightarrow$5pt.
\item \texttt{paper/tables/tab\_granular\_\{tissue,derm,oct,asym\}.tex} --- \verb|\tabcolsep|
      3$\rightarrow$8pt, widening $\sim$40\,\%-textwidth tables sitting under full-width captions.
\end{itemize}

\paragraph{Validation (second full pass).} Rebuilt \texttt{pdflatex} $\rightarrow$ \texttt{bibtex}
$\rightarrow$ \texttt{pdflatex}$\times$2, 31 pages, and re-rasterised every page. Former pages
28--31 (Tables 12--15) now sit \textbf{flush at the top} with the slack below --- visually confirmed
on pages 28, 30 and 31. Blue markup renders correctly (pages 9, 10). Zero overfull/underfull
\textbf{hboxes}. Seven \texttt{Underfull \textbackslash vbox \dots while \textbackslash output is
active} warnings appeared, all on the float pages: this is the \emph{expected signature of the fix}
--- the previous \verb|0pt plus 1fil| glue absorbed any slack with zero badness (precisely why the
floats were centred), and pinning to the top necessarily leaves a measurable gap at the bottom.
Cosmetically correct, and TMLR does not require \verb|\flushbottom| pages to be exactly full.

\begin{callout}{\textcolor{warnred}{One caveat}}
\texttt{paper/scripts/make\_granular\_tables.py} still emits \verb|\tabcolsep{3pt}|. It was left
unmodified (an attempted sync corrupted the escaping and was reverted). \textbf{Re-running that
generator will silently undo the four granular-table fixes.} Either edit the \texttt{.tex} files by
hand or fix the generator first. Consistent with \texttt{paper/data/README\_DATA.md}, which already
notes the committed tables are hand-maintained and their generators superseded.
\end{callout}

\paragraph{Build note.} \texttt{docs/main.tex} does not compile in \texttt{docs/}. Either build in
Overleaf (the professor's source of truth) or copy it to \texttt{paper/} ---
\texttt{paper/main\_rev.tex} / \texttt{main\_rev.pdf} is that working copy.

\end{document}
